% title:          Roots on the unit circle
% area:           polynomial-roots
% methods:        contradiction
% difficulty:
% difficulty-ai:  medium
% status:
% review:         none
% --- history: all optional, leave EMPTY if unknown ---
% origin:         putnam
% origin-date:    1989-12-02
% origin-number:  A3
% also-in:        book
% taken-from:
% references:     Putnam 1989, A3 (exam of 2 Dec 1989) - https://kskedlaya.org/putnam-archive/1989.pdf; Kalva archive (statement and z^9-modulus solution) - https://prase.cz/kalva/putnam/putn89.html and https://prase.cz/kalva/putnam/psoln/psol893.html; official report: Klosinski, Alexanderson, Larson, The Fiftieth William Lowell Putnam Mathematical Competition, Amer. Math. Monthly 98 (1991), no. 4, 319-327, doi:10.1080/00029890.1991.12000760 - https://api.crossref.org/works/10.1080/00029890.1991.12000760; Kedlaya, Poonen, Vakil, The William Lowell Putnam Mathematical Competition 1985-2000, MAA 2002, 1989 A3, p. 101 ff. (several solutions) - https://archive.org/details/Putnam19852000; Andreescu, Andrica, Complex Numbers from A to ... Z, 2nd ed., Birkhauser 2014, Sect. 5.2 'Algebraic Equations and Polynomials', Problem 7, p. 194 - https://archive.org/details/complexnumbersfr0000andr_b4d8 (bibliographic record https://api.crossref.org/works/10.1007/978-0-8176-8415-0); solutions sheet - https://www.sfu.ca/~lockhart/richard/Putnam1989Sol.pdf
% history:        ai
\begin{problem}
Prove that if the complex number \( z \) satisfies the equation

\[
11z^{10} + 10i z^9 + 10i z -11 = 0,
\]

then \( |z| = 1 \). (Here, \( i \) is the imaginary unit satisfying \( i^2 = -1 \).)
\end{problem}
\begin{solution}
We begin by expressing $z^9$ in terms of $z$:
\begin{equation*}
    z^9 = \frac{11 - 10iz}{11z + 10i}.
\end{equation*}
Taking the modulus on both sides, we examine the numerator and denominator separately:
\begin{align*}
    |11 - 10iz|^2 &= (11 - 10iz)(11 + 10iz) = 121 + 100|z|^2 + 220 \Im z, \\
    |11z + 10i|^2 &= (11z + 10i)(11\bar{z} - 10i) = 100 + 121|z|^2 + 220 \Im z.
\end{align*}
If we define $|z^9| = N/D$ where
\begin{equation*}
    N = 121 + 100|z|^2 + 220 \Im z, \quad D = 100 + 121|z|^2 + 220 \Im z,
\end{equation*}
then their difference is given by
\begin{equation*}
    N - D = 21(1 - |z|^2).
\end{equation*}

Now, consider the cases:
\\\\
- If $|z| > 1$, then $1 - |z|^2 < 0$, so $N < D$ and hence $|z^9| < 1$.
\\\\
- If $|z| < 1$, then $1 - |z|^2 > 0$, so $N > D$ and hence $|z^9| > 1$.

Both cases is a contradiction. Thus, we conclude that $|z| = 1$ for all roots $z$.
\end{solution}
